% ============================================================
\appendix
\section{S\l{}ownik notacji}\label{app:notacja}
% ============================================================

\begin{longtable}{@{}>{\raggedright}p{2.8cm}p{10cm}@{}}
\toprule
\textbf{Symbol} & \textbf{Znaczenie} \\
\midrule
\endfirsthead
\toprule
\textbf{Symbol} & \textbf{Znaczenie} \\
\midrule
\endhead
\bottomrule
\endlastfoot

\multicolumn{2}{@{}l}{\textbf{Fundamenty}} \\
\midrule
$\Nzero$ & Absolutna nico\'s\'c: brak przestrzeni, materii, czasu \\
$\Ffield$ & Funkcja stanu: miara odej\'scia od nico\'sci; $\Ffield = 0$ w~$\Nzero$, $\Ffield > 0$ w~$\mathcal{S}_1$ \\
$\mathcal{S}_0$ & Sektor nico\'sci: $\Ffield = 0$ \\
$\mathcal{S}_1$ & Sektor istnienia: $\Ffield > 0$ \\
$\Phi$ & Pole przestrzenno\'sci: g\k{e}sto\'s\'c wygenerowanej przestrzeni \\
$\PhiZero$ & Warto\'s\'c referencyjna $\Phi$ (warunki laboratoryjne) \\
\midrule

\multicolumn{2}{@{}l}{\textbf{\'Zr\'od\l{}a i~interferencja}} \\
\midrule
$q_a$ & Sta\l{}a generacji przestrzennej cz\k{a}stki typu~$a$;
  $[q] = [\mathrm{M}^{-1}\mathrm{L}]$.
  Metryka minimalna: $q_{\mathrm{min}} = 8\pi G_0/c_0^2$ (tw.~\ref{thm:newton});
  eksponencjalna: $q_{\mathrm{exp}} = 4\pi G_0/c_0^2$ (uw.~\ref{rem:metric-relation}) \\
$\rho_a$ & G\k{e}sto\'s\'c masy--energii \'zr\'od\l{}a typu~$a$ \\
$\phi(r)$ & Profil generacji: kszta\l{}t pola od pojedynczego \'zr\'od\l{}a \\
$\mathcal{D}[\Phi]$ & Operator dynamiki pola $\Phi$ \\
$\mathcal{N}[\Phi]$ & Cz\l{}on samointerferencji (nieliniowy):
  $\alpha(\nabla\Phi)^2/(\PhiZero\Phi)
  + \beta\Phi^2/\PhiZero^2
  - \gamma\Phi^3/\PhiZero^3$ \\
$\alpha$ & Bezwymiarowa sta\l{}a sprz\k{e}\.zenia gradientowego;
  $\alpha = 2$ z~wariacji dzia\l{}ania (wn.~\ref{cor:alpha-fixed}) \\
$\beta$ & Sta\l{}a odpychania $[\mathrm{L}^{-2}]$ (cz\l{}on kubiczny);
  bezwymiarowo: $\hat\beta = \beta r_0^2$ \\
$\gamma$ & Sta\l{}a studni $[\mathrm{L}^{-2}]$ (cz\l{}on kwartyczny);
  bezwymiarowo: $\hat\gamma = \gamma r_0^2$ \\
$E_{\mathrm{int}}(d)$ & Energia oddzia\l{}ywania dw\'och \'zr\'ode\l{} w~odleg\l{}o\'sci $d$ \\
$\varphi(t)$ & $\Phi(t)/\PhiZero$ --- bezwymiarowe pole kosmologiczne
  (to\.zsame z~$\psi$ w~kontek\'scie FRW);
  \textbf{uwaga:} w~sekcji ciemnej energii u\.zywane te\.z
  jako $\varphi = \Phi/\PhiZero - 1$ (perturbacja) ---
  znaczenie wynika z~kontekstu \\
$\psi$ & $\Phi/\PhiZero$ --- bezwymiarowe pole w~r\'ownaniach
  dynamicznych i~analizie stabilno\'sci;
  $\psi \equiv \chi$ (r\'o\.zne litery,
  ta sama wielko\'s\'c); $\varphi = \psi$ gdy u\.zyte
  jako pole kosmologiczne, $\varphi = \psi - 1$ gdy u\.zyte
  jako perturbacja \\
\midrule

\multicolumn{2}{@{}l}{\textbf{Trzy re\.zimy}} \\
\midrule
$r_{\mathrm{well}}$ & Skala studni przestrzennej (sklejenie) \\
$r_{\mathrm{rep}}$ & Skala odpychania przestrzennego \\
$V_{\mathrm{eff}}(r)$ & Efektywny potencja\l{} mi\k{e}dzy dwoma \'zr\'od\l{}ami \\
\midrule

\multicolumn{2}{@{}l}{\textbf{Dynamiczne sta\l{}e}} \\
\midrule
$c(\Phi)$ & Lokalna pr\k{e}dko\'s\'c \'swiat\l{}a: $c_0(\PhiZero/\Phi)^{1/2}$ \\
$\hbar(\Phi)$ & Lokalna sta\l{}a Plancka: $\hbar_0(\PhiZero/\Phi)^{1/2}$ \\
$G(\Phi)$ & Lokalna sta\l{}a grawitacyjna: $G_0(\PhiZero/\Phi)$ \\
$\lP$ & D\l{}ugo\'s\'c Plancka (sta\l{}a, niezmienna) \\
$c_0, \hbar_0, G_0$ & Obserwowane warto\'sci sta\l{}ych w~warunkach referencyjnych \\
\midrule

\multicolumn{2}{@{}l}{\textbf{Kosmologia i~czarne dziury}} \\
\midrule
$\Delta\Phi$ & Nadwy\.zka przestrzeni z~fluktuacji (\'zr\'od\l{}o ciemnej energii) \\
$w_{\mathrm{DE}}(z)$ & R\'ownanie stanu ciemnej energii \\
$\Srel$ & Entropia relacyjna: $-\int \rho_{\mathrm{cg}}\ln\rho_{\mathrm{cg}}\,\sqrt{h}\,d^3x$ \\
$g_{\mu\nu}^{\mathrm{eff}}$ & Metryka efektywna wynikaj\k{a}ca z~konfiguracji $\Phi$ \\
$\mathcal{I}(\mathcal{R})$ & Przepustowo\'s\'c informacyjna regionu $\mathcal{R}$ \\
\midrule

\multicolumn{2}{@{}l}{\textbf{Formalizm i~sp\'ojno\'s\'c}} \\
\midrule
$\Ffield_{\!\mathrm{abs}}(x)$ & Absolutna funkcja stanu: $\Phi(x)^2/\PhiZero^2$; $= 0$ w~$\Nzero$ (def.~\ref{def:Fabs}) \\
$\Floc(x)$ & Perturbacyjna funkcja stanu: $(\Phi(x) - \PhiZero)^2 = \PhiZero^2\varphi^2$ (def.~\ref{def:Floc}) \\
$\Delta(x)$ & Lokalna asymetria chiralna ze znakiem ($\langle\hat{s}\rangle$); $\int\Delta\sqrt{h}\,d^3x = 0$ (ZS1) \\
$T_{\mu\nu}^{(\Phi)}$ & Efektywny tensor energii-p\k{e}du pola $\Phi$ (stw.~\ref{prop:T-Phi}) \\
$\mathcal{T}$ & Zbi\'or konfiguracji trywialnych (brak struktury) \\
$\rho_{\mathrm{cg}}(x)$ & Grubozarniona g\k{e}sto\'s\'c substratu \\
$\mu_\beta$ & Miara Gibbsa substratu \\
$\beta_c$ & Krytyczna temperatura odwrotna substratu \\
$G_{\mathrm{eff}}(k,a)$ & Efektywna sta\l{}a grawitacyjna zale\.zna od skali \\
$\alpha_{\mathrm{eff}}$ & Bezwymiarowa sta\l{}a sprz\k{e}\.zenia: $q\PhiZero/(4\pi)$ \\
$m_{\mathrm{sp}}$ & Przestrzenna masa efektywna: $\sqrt{3\gamma - 2\beta}$ (stw.~\ref{prop:vacuum-stability});
  \textbf{dla $\beta = \gamma$}: $m_{\mathrm{sp}} = \sqrt{\gamma}$, a~zatem
  $m_{\mathrm{sp}}^2 = \gamma > 0$ (perturbacje stabilne) \\
$m_{\mathrm{cosmo}}$ & Kosmologiczna masa efektywna: $\sqrt{4c_0^2\gamma/3}$ z~r\'ownania~\eqref{eq:Phi-cosmo-linearized}; r\'o\.zna od $m_{\mathrm{sp}}$ (uwaga~\ref{rem:two-masses}) \\
$\tau_0$ & Kosmologiczna skala czasowa: $\sqrt{\PhiZero}/(c_0\sqrt{3\gamma - 2\beta})$; dla $\beta=\gamma$: $\sqrt{\PhiZero}/(c_0\sqrt{\gamma})$ (def.~\ref{def:tau0}) \\
$f(\Phi)$ & Zmodyfikowany cz\l{}on kinetyczny: $1 + 2\alpha\ln(\Phi/\PhiZero)$ \\
$j_{\Ffield}^\mu$ & Pr\k{a}d funkcji stanu: $2(\Phi-\PhiZero)\partial^\mu\Phi$ (stw.~\ref{prop:F-current}) \\
$\chi$ & Bezwymiarowe pole $\Phi/\PhiZero$; to\.zsame z~$\psi$
  (u\.zywane g\l{}\'ownie w~sek.~\ref{sec:czarne}
  i~dod.~\ref{app:ringdown}) \\
$\iota$ & Inwersja dualno\'sci $\Nzero$: $\chi \mapsto 1/\chi$ (tw.~\ref{thm:N0-duality}) \\
\midrule

\multicolumn{2}{@{}l}{\textbf{Defekty i~materia}} \\
\midrule
$m_{\mathrm{def}}$ & Masa radialnego kinka (defektu $\Phi$): $4\pi\int[\frac{1}{2}(\Phi')^2 + V_{\mathrm{eff}} - V_{\mathrm{eff}}(\PhiZero)]\,r^2 dr$ \\
$\Phi_{\!0}$ & Warto\'s\'c pola $\Phi$ w~centrum defektu ($\neq \PhiZero$) \\
$n$ & Liczba w\k{e}z\l{}\'ow profilu radialnego --- pokolenie cz\k{a}stki \\
$e^a{}_\mu$ & Tetrada TGP: $g_{\mu\nu}^{\mathrm{eff}} = \eta_{ab}\,e^a{}_\mu\,e^b{}_\nu$ (def.~\ref{def:tetrada}) \\
$\omega_\mu{}^{bc}$ & Koneksja spinowa wyznaczona z~tetrady TGP \\
$\Psi$ & Pole Diraca sprz\k{e}\.zone z~$g_{\mu\nu}^{\mathrm{eff}}(\Phi)$ (hip.~\ref{hyp:dirac-TGP}) \\
$n_Q$ & \L{}adunek topologiczny defektu (hip.~\ref{hyp:charge-topology}); $n_Q \in \mathbb{Z}$ \\
$A_\mu^a$ & Pola cechowania emergentne z~substratu (hip.~\ref{hyp:gauge-substrate}) \\
\midrule

\multicolumn{2}{@{}l}{\textbf{Ciemna energia --- formalizm}} \\
\midrule
$U(\varphi)$ & Potencja\l{} samosprz\k{e}\.zenia;
  $\varphi = \Phi/\PhiZero - 1$ w~sek.~\ref{sec:ciemna},
  $\varphi = \Phi/\PhiZero$ (tj.~$\psi$) w~sek.~8 ---
  konwencja wynika z~kontekstu \\
$\Lambda_{\mathrm{eff}}$ & Efektywna sta\l{}a kosmologiczna: resztkowy potencja\l{} pr\'o\.zniowy \\
$\varphi_{\min}$ & Po\l{}o\.zenie minimum $U(\varphi)$; $\neq 0$ z~zasady zerowej sumy \\
\midrule

\multicolumn{2}{@{}l}{\textbf{Topologia i~domeny}} \\
\midrule
$\mathcal{M}_k$ & Domena metryczna: sp\'ojny region z~$\Phi > 0$, $c > 0$ \\
\midrule

\multicolumn{2}{@{}l}{\textbf{Sektor nielokalny}} \\
\midrule
$K(x{-}y)$ & J\k{a}dro nielokalne z~korelator\'ow substratu \\
$\xi_{\mathrm{corr}}$ & D\l{}ugo\'s\'c korelacji substratu \\
$\eta$ & Wyk\l{}adnik anomalnego skalowania \\
$G_\varphi(k)$ & Propagator perturbacji $\delta\Phi$: $1/(k^2 + m_{\mathrm{eff}}^2)$ \\
$k_{\max}$ & Obci\k{e}cie UV z~substratu: $\sim 1/\lP$ \\
\midrule

\multicolumn{2}{@{}l}{\textbf{Substrat --- parametry}} \\
\midrule
$\mu$ & Masa kinetyczna operatora $\hat{s}_i$ \\
$m_0^2$ & Kwadrat masy go\l{}ej w~potencjale substratowym \\
$\lambda_0$ & Sta\l{}a samosprz\k{e}\.zenia $\hat{s}^4$ \\
$J$ & Sta\l{}a sprz\k{e}\.zenia mi\k{e}dzy s\k{a}siadami \\
$z$ & Koordynacja substratu (liczba s\k{a}siad\'ow) \\
$v$ & Pr\'o\.zniowa warto\'s\'c oczekiwana: $v^2 = (Jz - m_0^2)/\lambda_0$ \\
$C_\beta, C_\gamma$ & Wsp\'o\l{}czynniki coarse-grainingu: $\beta,\gamma = f(\lambda_0, v, C_{\beta,\gamma})$ \\
\end{longtable}

% ----------------------------------------------------------
\subsection{Tablica konwersji wymiarowej}\label{app:A-wymiary}
% ----------------------------------------------------------

W~r\'o\.znych cz\k{e}\'sciach tekstu sta\l{}e $\beta$, $\gamma$ wyst\k{e}puj\k{a}
zar\'owno z~wymiarem $[\mathrm{L}^{-2}]$ (uk\l{}ad fizyczny) jak i~jako bezwymiarowe
$\hat\beta, \hat\gamma$ (uk\l{}ad siatki). Poni\.zsza tablica zbiera konwencje.

\medskip\noindent
\textbf{A.~Wielko\'sci fundamentalne:}

\begin{center}
\begin{tabular}{@{}llll@{}}
\toprule
\textbf{Symbol} & \textbf{Wymiar} & \textbf{Definicja / skalowanie} & \textbf{Uwagi} \\
\midrule
$\Phi$ & $[-]$ & pole skalarne (bezwymiarowe) & generowane przez \'zr\'od\l{}a \\
$\PhiZero$ & $[-]$ & warto\'s\'c referencyjna $\Phi$ & ,,pr\'o\.znia'' ($\Phi\to\PhiZero$ w~$r\to\infty$) \\
$q$ & $[-]$ & si\l{}a \'zr\'od\l{}owania & $\mathcal D[\Phi_a]=q_a\rho_a$ \\
$\rho$ & $[\mathrm{L}^{-3}]$ & g\k{e}sto\'s\'c \'zr\'od\l{}owa & $\int\rho\,d^3x = 1$ (normalizacja) \\
$c(\Phi)$ & $[\mathrm{LT}^{-1}]$ & $c_0(\PhiZero/\Phi)^{1/2}$ & aks.~A6 \\
$\hbar(\Phi)$ & $[\mathrm{ML}^2\mathrm{T}^{-1}]$ & $\hbar_0(\PhiZero/\Phi)^{1/2}$ & aks.~A7 \\
$G(\Phi)$ & $[\mathrm{M}^{-1}\mathrm{L}^3\mathrm{T}^{-2}]$ & $G_0(\PhiZero/\Phi)$ & aks.~A8 \\
$\lP$ & $[\mathrm{L}]$ & $\sqrt{\hbar G/c^3} = \mathrm{const}$ & \textbf{jedyny niezmiennik} \\
\bottomrule
\end{tabular}
\end{center}

\medskip\noindent
\textbf{B.~Parametry r\'ownania pola i~substratu:}

\begin{center}
\begin{tabular}{@{}lllll@{}}
\toprule
\textbf{Symbol} & \textbf{Wymiar fiz.} & \textbf{Bezwym.} & \textbf{Konwersja} & \textbf{U\.zycie} \\
\midrule
$\beta$ & $[\mathrm{L}^{-2}]$ & $\hat\beta = \beta\,r_0^2$ & $\beta = \hat\beta/r_0^2$ & r\'ow.~pola~\eqref{eq:full-field-alt} \\
$\gamma$ & $[\mathrm{L}^{-2}]$ & $\hat\gamma = \gamma\,r_0^2$ & $\gamma = \hat\gamma/r_0^2$ & potencja\l{}, masa \\
$a_\Gamma$ & $[\mathrm{L}]$ & $1$ (jednostka siatki) & $r = a_\Gamma \xi$ & skrypty, solitony \\
$m_{\mathrm{sp}}$ & $[\mathrm{L}^{-1}]$ & $\hat m = m_{\mathrm{sp}} a_\Gamma$ & $m_{\mathrm{sp}} = \sqrt{\gamma}$ & profil Yukawa \\
$\kappa$ & $[-]$ & $\kappa = 3/(4\PhiZero)$ & --- & sek.~\ref{ssec:unified-action} \\
$K_{\rm geo}$ & $[-]$ & $K_{\rm geo} = 1$ (normalizacja) & --- & sprz.~kinetyczne \\
$V(\psi)$ & $[\mathrm{L}^{-2}]$ & $\hat V = V/\gamma$ & $V = \gamma\hat V$ & potencja\l{}~\eqref{eq:V-orig} \\
$g_0$ & $[-]$ & $g_0 = \Phi_0/\PhiZero$ & --- & amplituda solitonu \\
$A_{\rm tail}$ & $[-]$ & --- & --- & ogon oscylacyjny \\
$r_{21}$ & $[-]$ & $r_{21} = (A_{\rm tail}^\mu/A_{\rm tail}^e)^4$ & --- & stosunek mas \\
$\alpha_K$ & $[-]$ & $\alpha_K = 2\alpha\gamma/(\gamma r_0^2)$ & patrz dod.~\ref{app:ogon-masy} & sprz.~soliton. \\
\bottomrule
\end{tabular}
\end{center}

\medskip\noindent
\textbf{Konwencja skrypt\'ow:} Wszystkie skrypty w~\texttt{scripts/}
u\.zywaj\k{a} uk\l{}adu bezwymiarowego z~$a_\Gamma = 1$ (jednostka siatki)
i~$\gamma = a_\Gamma = 0{,}040049$ (fizyczna warto\'s\'c). Przej\'scie
do uk\l{}adu fizycznego: $r_{\rm phys} = a_\Gamma\,\xi$,
$\Phi_{\rm phys} = \PhiZero\,\psi$.

% ----------------------------------------------------------
\subsection{Aksjomaty TGP --- podsumowanie}

\begin{center}
\begin{tabular}{@{}cp{12.5cm}@{}}
\toprule
\textbf{Nr} & \textbf{Tre\'s\'c} \\
\midrule
A1 & Przestrze\'n wynika z~obecno\'sci materii
     (aksjomat~\ref{ax:przestrzen}) \\
A2 & $\Nzero$ jest stanem odniesienia z~$\Ffield = 0$, $\Phi = 0$
     (aksjomat~\ref{ax:N0}); wyprowadzalny z~$\Ffield$
     (stw.~\ref{prop:N0-from-F}) \\
A3 & Zerowa suma: dwa niezale\.zne warunki ---
     (ZS1) chiralna: $\int\Delta\sqrt{h}\,d^3x = 0$;
     (ZS2) przestrzenna: $\int(\Phi-\PhiZero)\sqrt{h}\,d^3x = 0$
     (aksjomat~\ref{ax:zero}, precyzacja~\ref{rem:zero-precyzacja}) \\
A4 & $\Nzero$ jest niestabilny: $\mu(\Ffield=0) = 0$
     (aksjomat~\ref{ax:P3}); wyprowadzalny z~termodynamiki
     substratu (stw.~\ref{prop:P3-deriv}) i~z~potencja\l{}u
     (stw.~\ref{prop:P3-action}) \\
A5 & Ka\.zde \'zr\'od\l{}o materii generuje wk\l{}ad do $\Phi$
     (aksjomat~\ref{ax:zrodlo}) \\
A6 & $c(\Phi) = c_0(\PhiZero/\Phi)^{1/2}$
     (aksjomat~\ref{ax:c}); wyprowadzalny z~metryki konforemnej
     (stw.~\ref{prop:conformal-unique}) \\
A7 & $\hbar(\Phi) = \hbar_0(\PhiZero/\Phi)^{1/2}$
     (aksjomat~\ref{ax:hbar}); \textbf{konsekwencja} A6 + sta\l{}o\'s\'c
     $\lP$ + sp\'ojno\'s\'c propagatora (tw.~\ref{thm:exponents}) \\
A8 & $G(\Phi) = G_0(\PhiZero/\Phi)$
     (aksjomat~\ref{ax:G}); \textbf{konsekwencja} A6\,+\,A7 + sta\l{}o\'s\'c
     $\lP$ (tw.~\ref{thm:exponents}) \\
\bottomrule
\end{tabular}
\end{center}

% ----------------------------------------------------------
\subsection{Za\l{}o\.zenia bazowe N-body TGP (N0-1 do N0-7)}
\label{ssec:N0-lista}

Poni\.zsza tabela definiuje \textbf{kanoniczne za\l{}o\.zenia N0-1 do N0-7}
u\.zywane w~obliczeniach N-body i~w~sektorze FDM.
S\k{a} one logicznie zale\.zne od aksjomat\'ow A1--A8 i~mo\.zna je traktowa\'c
jako ich \textbf{graniczne (quasi-klasyczne) konsekwencje} przy $\Phi \approx \PhiZero$.

\begin{center}
\begin{tabular}{@{}clp{7.5cm}l@{}}
\toprule
\textbf{Nr} & \textbf{Tre\'s\'c}
  & \textbf{Jawna posta\'c} & \textbf{\'Zr\'od\l{}o} \\
\midrule
N0-1
  & Pole $\Phi$ jako jedyne \'zr\'od\l{}o metryki
  & $g_{\mu\nu} = g_{\mu\nu}[\Phi]$;
    $\Phi(x) = \PhiZero + \delta\Phi(x)$,
    $|\delta\Phi| \ll \PhiZero$
  & A1, A5 \\[4pt]
N0-2
  & T\l{}o referencyjne: jednorodne pr\'o\.zniowe $\Phi$
  & $\Phi_{\mathrm{bg}} = \PhiZero = \mathrm{const}$;
    warunek pr\'o\.zniowy $\beta = \gamma$
    (stw.~\ref{prop:vacuum-condition})
  & A1, thm.~\ref{prop:vacuum-condition} \\[4pt]
N0-3
  & Profil Yukawy dla pojedynczego \'zr\'od\l{}a
  & $\delta\Phi(r) = -C\,e^{-m_{\mathrm{sp}}\,r}/r$;
    $C > 0$ --- amplituda generacji przestrzennej
  & A5, stw.~\ref{prop:vacuum-stability} \\[4pt]
N0-4
  & Zwi\k{a}zek amplitudy z~mas\k{a} \'zr\'od\l{}a
    (\textbf{Droga B})
  & $C = q\,M\,\PhiZero / (4\pi)$;
    w~granicy izotropowej: $C = m_{\mathrm{sp}}/(2\sqrt{\pi})$
    (skrypt \texttt{ex39})
  & A5, stw.~\ref{prop:composite-profile} \\[4pt]
N0-5
  & Warunek pr\'o\.zniowy i~posta\'c potencja\l{}u efektywnego
  & $\beta = \gamma$;
    $V_{\mathrm{eff}}(g) = \frac{\gamma}{3}g^3 - \frac{\gamma}{4}g^4$\footnote{Formulacja kanoniczna TGP u\.zywa potencja\l{}u $P(g) = (\beta/7)\,g^7 - (\gamma/8)\,g^8$ z~$K(g) = g^4$; powy\.zsza forma~$(3,4)$ jest wczesn\k{a} wersj\k{a} dziel\k{a}c\k{a} t\k{e} sam\k{a} RHS solitonu.},
    gdzie $g = \Phi/\PhiZero$
    (stw.~\ref{prop:vacuum-condition}, \ref{prop:nonlinear-stability})
  & A1, aksjomat struktury $\Gamma$ \\[4pt]
N0-6
  & Masa pola przestrzenno\'sci (masa skalaru TGP)
  & $m_{\mathrm{sp}}^2 = 3\gamma - 2\beta = \gamma$
    (dla $\beta = \gamma$);
    $m_{\mathrm{sp}} = \sqrt{\gamma}$
    (stw.~\ref{prop:vacuum-stability}, prop.~\ref{prop:eps-th})
  & N0-5 \\[4pt]
N0-7
  & Zmienno\'s\'c $G$ i~ograniczenie kosmologiczne
  & $G(\Phi) = G_0\,\PhiZero/\Phi$ (A8);
    przy $\Phi_{\mathrm{bg}} = \PhiZero\,\psi_{\mathrm{bg}}(t)$:
    $\dot{G}/G = -\dot\psi/\psi$.
    Dla zamro\.zonego pola ($|\dot\psi| \ll H\psi$, stw.~\ref{prop:Lambda-eff}):
    $|\dot{G}/G| \lesssim H_0\,\delta_\psi$,
    gdzie $\delta_\psi = |\psi_{\rm ini} - 1|\cdot e^{-\xi}$,
    $\xi$ --- ca\l{}kowite t\l{}umienie Hubble'a
    od momentu pocz\k{a}tkowego do $t_0$.
    Przy $\psi_{\rm ini} = 7/6$:
    (a)~oszacowanie analityczne ($\xi \approx 3/2$):
    $\delta_\psi \approx 0{,}037$;
    (b)~z~dominacj\k{a} DE ($\xi \approx 2$):
    $\delta_\psi \approx 0{,}023$;
    (c)~\textbf{pe\l{}na ewolucja FRW} (skrypt \texttt{p\_frw\_full\_evolution.py},
    Radau ODE, $z = 10^9 \to 0$):
    \textbf{Historycznie} przy starym $\kappa_{\rm old} = 3/(2\Phi_0) = 0{,}061$:
    $|\dot{G}/G|/H_0 = 0{,}035$ --- czynnik $1{,}75\times$ powy\.zej LLR.
    \textbf{Po korekcji} elementu obj\k{e}to\'sciowego
    $\sqrt{-g} = c_0\psi$ (prop.~\ref{prop:kappa-corrected},
    sek.~\ref{ssec:unified-action}):
    \begin{equation*}
      \kappa = \frac{3}{4\Phi_{\rm eff}} = \frac{7}{2\PhiZero} \approx 0{,}030,
    \end{equation*}
    co daje $|\dot{G}/G|/H_0 \lesssim 0{,}009$--$0{,}02$
    --- \textbf{zgodne z~LLR} (Williams et~al.\ 2004:
    $|\dot{G}/G| < 0{,}02\,H_0$).
    \textbf{Status: ZAMKNI\k{E}TY} (napi\k{e}cie zlikwidowane
    poprawnym $\kappa$)
  & A8, N0-5, stw.~\ref{prop:psi-ini-derived} \\
\bottomrule
\end{tabular}
\end{center}

\begin{remark}[Hierarchia aksjomat\'ow i~za\l{}o\.ze\'n N0]
Za\l{}o\.zenia N0-1 do N0-6 s\k{a} \textbf{logi\-cznie zale\.zne}
od aksjomat\'ow A1--A8 i~stanowi\k{a} ich spe\-cjali\-zacj\k{e}
do re\.zimu quasi-klasycznego ($\Phi \approx \PhiZero$, s\l{}abe pole).
Kolejno\'s\'c implikacji:
\begin{equation*}
  \underbrace{A1}_{\text{mat.}\to\Phi}
  \;\Rightarrow\;
  \underbrace{A5}_{\text{\'zr\'od\l{}o}}
  \;\Rightarrow\;
  \underbrace{\text{N0-1}}_{\delta\Phi\ll\PhiZero}
  \;\Rightarrow\;
  \underbrace{\text{N0-2}}_{\beta=\gamma}
  \;\Rightarrow\;
  \underbrace{\text{N0-3}}_{\text{Yukawa}}
  \;\Rightarrow\;
  \underbrace{\text{N0-4}}_{C=qM\Phi_0/(4\pi)}
  \;\Rightarrow\;
  \underbrace{\text{N0-5}}_{V_{\mathrm{eff}}}
  \;\Rightarrow\;
  \underbrace{\text{N0-6}}_{m_{\mathrm{sp}}=\sqrt{\gamma}}.
\end{equation*}
Ka\.zde N0-x mo\.zna traktowa\'c jako aksjomat roboczy \textbf{lub} jako
twierdzenie wyprowadzalne z~poprzednich --- zale\.znie od poziomu opisu.
\end{remark}

% -------------------------------------------------------------------
\subsection{Wyprowadzenie N0-5 $\Rightarrow$ N0-6: masa pola~$m_{\rm sp}$}%
\label{ssec:N0-derivation}
% -------------------------------------------------------------------

\begin{proposition}[N0-6 wynika z~N0-5]%
\label{prop:N0-6-from-N0-5}
Przy warunku pr\'o\.zniowym $\beta = \gamma$ (N0-5),
efektywny potencja\l{}
\begin{equation}\label{eq:Veff-N0-5}
  V_{\rm eff}(g) = \frac{\gamma}{3}\,g^3 - \frac{\gamma}{4}\,g^4,
  \qquad g = \frac{\Phi}{\PhiZero},
\end{equation}%
\footnote{Formulacja kanoniczna TGP u\.zywa potencja\l{}u $P(g) = (\beta/7)\,g^7 - (\gamma/8)\,g^8$ z~$K(g) = g^4$; powy\.zsza forma~$(3,4)$ jest wczesn\k{a} wersj\k{a} dziel\k{a}c\k{a} t\k{e} sam\k{a} RHS solitonu.}%
ma minimum w~$g_{\rm min} = 1$ (pr\'o\.znia) i~mas\k{e} pola:
\begin{equation}\label{eq:msp-N0-6}
  \boxed{m_{\rm sp}^2 = V_{\rm eff}''(1) = 3\gamma - 2\beta \;\overset{\beta=\gamma}{=}\; \gamma,
  \qquad m_{\rm sp} = \sqrt{\gamma}.}
\end{equation}
\end{proposition}

\begin{proof}
\emph{Krok~1 (minimum).}
$V_{\rm eff}'(g) = \gamma g^2 - \gamma g^3 = \gamma g^2(1-g)$.
Zera: $g = 0$ (niestabilne) i~$g = 1$ (fizyczne minimum).

\emph{Krok~2 (druga pochodna przy $g=1$).}
$V_{\rm eff}''(g) = 2\gamma g - 3\gamma g^2$.
Przy $g = 1$:
$V_{\rm eff}''(1) = 2\gamma - 3\gamma = -\gamma$.

\noindent\textbf{Uwaga o~znaku:} $V_{\rm eff}''(1) = -\gamma < 0$
sugeruje \emph{nie minimum}. Wyja\'snienie: $V_{\rm eff}$
jest \textbf{potencja\l{}em rolorodnym} (t\l{}umienie, nie oscylacja).
Stabilno\'s\'c pr\'o\.zni sprawdzamy przez liniow\k{a}
perturbacj\k{e} r\'ownania pola, nie przez drugą pochodn\k{a}~$V$.

\emph{Krok~3 (linearyzacja r\'ownania pola).}
Pe\l{}ne r\'ownanie pola TGP dla $g = 1 + \delta g$, $|\delta g| \ll 1$:
\begin{equation}\label{eq:field-linearized}
  \frac{1}{\PhiZero}\nabla^2(\delta g)
  + \frac{2(\nabla\delta g)^2}{\delta g}
  + V_{\rm eff}'(1+\delta g) = 0.
\end{equation}
Linearyzacja ($(\nabla\delta g)^2 \to 0$, $V_{\rm eff}'(1+\delta g) \approx
V_{\rm eff}''(1)\,\delta g$):
\begin{equation}\label{eq:delta-g-linear}
  \nabla^2(\delta g) - m_{\rm sp}^2\,\delta g = 0,
  \qquad
  m_{\rm sp}^2 \;=\; -\PhiZero\,V_{\rm eff}''(1) \;=\; -\PhiZero(-\gamma)
  \;=\; \gamma\,\PhiZero.
\end{equation}
W~jednostkach $\PhiZero = 1$ (bezwymiarowe): $m_{\rm sp}^2 = \gamma$.

\emph{Krok~4 (profil Yukawy --- N0-3).}
R\'ownanie~\eqref{eq:delta-g-linear} ma sferyczno-symetryczne
rozwi\k{a}zanie zanikaj\k{a}ce w~niesko\'nczono\'sci:
$\delta g(r) = -\,C\,e^{-m_{\rm sp}\,r}/r$,
co jest dok\l{}adnie profilem Yukawy N0-3 z~\textbf{mas\k{a}
$m_{\rm sp} = \sqrt{\gamma}$}~(N0-6). \qed

\medskip\noindent
{\small \textbf{Weryfikacja numeryczna:} skrypt
\texttt{n0\_axioms\_chain\_verification.py}
(katalog \texttt{scripts/}), testy T-N0-5 i~T-N0-6: 2/2~PASS.}
\end{proof}

\begin{remark}[Interpretacja fizyczna $m_{\rm sp}$]\label{rem:msp-physics}
Masa $m_{\rm sp} = \sqrt{\gamma}$ okre\'sla \textbf{zasi\k{e}g}
odchyle\'n od pr\'o\.zni: perturbacja $\delta\Phi$ na skali
$r \ll m_{\rm sp}^{-1}$ jest silna (re\.zim wewn\k{e}trzny),
a~dla $r \gg m_{\rm sp}^{-1}$ zanika eksponencjalnie (re\.zim Yukawa).

W~kosmologii: $m_{\rm sp}^{-1} = 1/\sqrt{\gamma}$.
Z~$\gamma \sim H_0^2/c_0^2$: $m_{\rm sp}^{-1} \sim c_0/H_0 = R_H$
(horyzont Hubble'a). Pole TGP ma zasi\k{e}g rz\k{e}du rozmiaru
Wszech\'swiata --- co jest sp\'ojne z~rol\k{a} $\Phi$ jako
pola generuj\k{a}cego przestrze\'n w~skali kosmologicznej.

Predykcja falsyfikowalna (Kill-shot K-G): cz\k{e}stotliwo\'s\'c
odci\k{e}cia modu oddechowego GW:
$f_{\rm cut} = c_0 m_{\rm sp}/(2\pi) \approx H_0/(2\pi)
\approx 3{,}7 \times 10^{-18}\,\mathrm{Hz}$.
\end{remark}

% -------------------------------------------------------------------
\subsection*{Disambiguacja trzech obiekt\'ow zwanych ,,potencja\l{}em''}
\label{app:A-potential-disambiguation}
% -------------------------------------------------------------------

W~tek\'scie TGP wyst\k{e}puj\k{a} trzy r\'o\.zne obiekty okre\'slane
s\l{}owem \emph{potencja\l{}}, pe\l{}ni\k{a}ce zasadniczo inne role.
Poni\.zsza tabela rozr\'o\.znia je jednoznacznie.

\begin{center}
\small
\begin{tabular}{@{}p{2.6cm}p{2.6cm}p{3.0cm}p{4.8cm}@{}}
\toprule
\textbf{Symbol} & \textbf{Nazwa} & \textbf{Argument} & \textbf{Rola i~lokalizacja} \\
\midrule
$V_{\mathrm{eff}}(r)$
  & Potencja\l{} mi\k{e}dzycz\k{a}stkowy
  & $r$ --- odleg\l{}o\'s\'c mi\k{e}dzy cz\k{a}stkami
  & Si\l{}a mi\k{e}dzycz\k{a}stkowa; trzy re\.zimy TGP;
    $F = -\partial_r V_{\mathrm{eff}}$;
    \textbf{sek03}, tw.~\ref{thm:uniqueness} \\
\midrule
$V_{\mathrm{mod}}(\varphi)$
  & Potencja\l{} samopola
  & $\varphi = \Phi/\PhiZero$ --- znormalizowane pole t\l{}a
  & Dynamika jednorodnego t\l{}a $\Phi$; decyduje o~ewolucji
    kosmologicznej i~stabilno\'sci pr\'o\.zni;
    $V_{\mathrm{mod}} = \beta\varphi^2 - \gamma\varphi^3$;
    \textbf{sek05/sek08} (\S\ref{ssec:rownanko}) \\
\midrule
$V_{\mathrm{pot}}(\psi)$
  & Potencja\l{} dzia\l{}ania kosmologicznego
  & $\psi = \sqrt{\Phi}$ --- pole pierwiastkowe
  & Dzia\l{}anie skalarne w~FLRW; powi\k{a}zany z~$w_{\mathrm{DE}}$;
    pojawia si\k{e} w~r\'ownaniu Kleina--Gordona dla $\psi$;
    \textbf{sek05/sek08} (\S\ref{ssec:architektura}) \\
\bottomrule
\end{tabular}
\end{center}

\begin{remark}[Zwi\k{a}zek mi\k{e}dzy potencja\l{}ami]\label{rem:A-potential-rel}
  Trzy potencja\l{}y s\k{a} ze sob\k{a} powi\k{a}zane, lecz \textbf{nie
  s\k{a} to\.zsame}:
  \begin{itemize}[nosep]
    \item $V_{\mathrm{eff}}(r)$ wynika z~ca\l{}kowania $\mathcal{N}[\Phi]$
      wzgl\k{e}dem po\l{}o\.zenia --- to potencja\l{} dwucz\k{a}stkowy,
      zale\.zny od geometrii (odleg\l{}o\'sci).
    \item $V_{\mathrm{mod}}(\varphi)$ to potencja\l{} jednocz\k{a}stkowy
      w~przestrzeni warto\'sci pola --- odpowiada czci cz\l{}onu
      $\mathcal{N}[\Phi]$ bez gradient\'ow ($\nabla\Phi = 0$).
    \item $V_{\mathrm{pot}}(\psi)$ jest reparametryzacj\k{a}
      $V_{\mathrm{mod}}$ dla pola $\psi = \sqrt\Phi$, wygodniejsz\k{a}
      w~zapisie kowariantnego dzia\l{}ania.
  \end{itemize}
  Gdy piszemy po prostu ,,potencja\l{} TGP'' bez dodatkowego
  okre\'slenia, zawsze powinno by\'c jasne z~kontekstu,
  kt\'ory z~trzech obiekt\'ow mamy na my\'sli.
\end{remark}

% -------------------------------------------------------------------
\subsection{Mapa zale\.zno\'sci logicznych TGP}\label{ssec:dependency-map}
% -------------------------------------------------------------------

Poni\.zsza mapa pokazuje \textbf{kluczowe \'la\'ncuchy dedukcji} od
aksjomat\'ow do predykcji testowalnych.
Dla ka\.zdego kroku podajemy: (i)~wej\'scie logiczne,
(ii)~wynik po\'sredni lub ko\'ncowy z~etykiet\k{a} statusu
(\S\ref{rem:status-levels}), (iii)~faz\k{e} rozwoju teorii.

\subsubsection*{\'La\'ncuch I\@.\ --- Substrat $\Gamma$ i~operator $\mathcal{D}$}
\begin{center}
\small
\begin{tabular}{@{}>{\raggedright}p{3.5cm}c>{\raggedright}p{4.5cm}cl@{}}
\toprule
\textbf{Wej\'scie} & $\Rightarrow$ & \textbf{Wynik} & \textbf{Status} & \textbf{Faza} \\
\midrule
A1 (substrat $\Gamma$, $K(\varphi)=\varphi^4$)
  & $\Rightarrow$
  & \texttt{prop:substrate-action}: dzia\l{}anie $S_\Gamma$
  & \statuslabel{Propozycja}
  & I.A \\
\texttt{prop:substrate-action}
  & $\Rightarrow$
  & $\alpha = 2$ (wn.~\ref{cor:alpha-fixed}): cz\l{}on gradientowy
  & \statuslabel{Twierdzenie}
  & I.A \\
\texttt{prop:substrate-action}
  & $\Rightarrow$
  & \texttt{thm:D-uniqueness}: jedyno\'s\'c operatora $\mathcal{D}$
  & \statuslabel{Twierdzenie}
  & I.B \\
$K(\varphi)=\varphi^4$, $\mathbb{Z}_2$
  & $\Rightarrow$
  & \texttt{rem:minimal-nonlinearity}: minimalno\'s\'c
  & \statuslabel{Propozycja}
  & I.A \\
A1, $K(\varphi)$, $\beta=\gamma$
  & $\Rightarrow$
  & \texttt{prop:vacuum-condition}: pr\'o\.znia stabilna
  & \statuslabel{Twierdzenie}
  & I.B \\
\texttt{prop:vacuum-condition}
  & $\Rightarrow$
  & \texttt{prop:vacuum-stability}: $m_{\mathrm{sp}}=\sqrt{\gamma}$
  & \statuslabel{Twierdzenie}
  & I.B \\
\bottomrule
\end{tabular}
\end{center}

\subsubsection*{\'La\'ncuch II\@.\ --- Sta\l{}e dynamiczne i~falsyfikowalno\'s\'c}
\begin{center}
\small
\begin{tabular}{@{}>{\raggedright}p{3.5cm}c>{\raggedright}p{4.5cm}cl@{}}
\toprule
\textbf{Wej\'scie} & $\Rightarrow$ & \textbf{Wynik} & \textbf{Status} & \textbf{Faza} \\
\midrule
A6 + sta\l{}o\'s\'c $\lP$
  & $\Rightarrow$
  & A7, A8: $\hbar(\Phi)$, $G(\Phi)$
  & \statuslabel{Twierdzenie}
  & I.C \\
A1--A8
  & $\Rightarrow$
  & \texttt{thm:exponents}: wyk\l{}adniki $(-\tfrac{1}{2}, -\tfrac{1}{2}, -1)$
  & \statuslabel{Twierdzenie}
  & I.C \\
A6 + jednorodne t\l{}o
  & $\Rightarrow$
  & $c_T = c_0$ (stw.~\ref{prop:cT}); zgodne z~GW170817
  & \statuslabel{Twierdzenie}
  & I.C \\
A1, A5, N0-1 do N0-4
  & $\Rightarrow$
  & \texttt{thm:newton}: grawitacja Newtona
  & \statuslabel{Twierdzenie}
  & II.A \\
\texttt{thm:newton} + PPN
  & $\Rightarrow$
  & $\gamma_{\mathrm{PPN}}=1$, $\beta_{\mathrm{PPN}}=1$
  & \statuslabel{Propozycja}
  & II.A \\
A8 + ewolucja $\Phi_{\mathrm{bg}}(z)$
  & $\Rightarrow$
  & $\Delta G/G \lesssim H_0\varepsilon$ (N0-7)
  & \statuslabel{Program}
  & II.B \\
\bottomrule
\end{tabular}
\end{center}

\subsubsection*{\'La\'ncuch III\@.\ --- Kosmologia i~perturbacje}
\begin{center}
\small
\begin{tabular}{@{}>{\raggedright}p{3.5cm}c>{\raggedright}p{4.5cm}cl@{}}
\toprule
\textbf{Wej\'scie} & $\Rightarrow$ & \textbf{Wynik} & \textbf{Status} & \textbf{Faza} \\
\midrule
A1, A3 (ZS2), $\beta=\gamma$
  & $\Rightarrow$
  & $\Lambda_{\mathrm{eff}} \approx c_0^2\gamma/3$ (stw.~\ref{prop:cosmo-system})
  & \statuslabel{Propozycja}
  & II.C \\
\texttt{prop:cosmo-system}
  & $\Rightarrow$
  & \texttt{sssec:cosmo-complete}: pe\l{}ny uk\l{}ad FRW
  & \statuslabel{Propozycja}
  & II.C \\
\texttt{prop:cosmo-system} + $\psi_{\mathrm{ini}}$
  & $\Rightarrow$
  & \texttt{prop:ghost-free-MS}: brak ghost\'ow, $c_s^2=c_0^2$
  & \statuslabel{Propozycja}
  & III.A \\
\texttt{prop:ghost-free-MS} + $v_k$ (num.)
  & $\Rightarrow$
  & $n_s - 1$, $r = P_T/P_S$ (C.2B --- \statuslabel{Program})
  & \statuslabel{Program}
  & III.B \\
T\l{}o kosmol.\ + $T_\Gamma(z)$
  & $\Rightarrow$
  & $w_{\mathrm{DE}}(z)$, $\eta_{\mathrm{slip}}$ (predykcje)
  & \statuslabel{Hipoteza}
  & II.C \\
\bottomrule
\end{tabular}
\end{center}

\subsubsection*{\'La\'ncuch IV\@.\ --- Materia, defekty i~generacje}
\begin{center}
\small
\begin{tabular}{@{}>{\raggedright}p{3.5cm}c>{\raggedright}p{4.5cm}cl@{}}
\toprule
\textbf{Wej\'scie} & $\Rightarrow$ & \textbf{Wynik} & \textbf{Status} & \textbf{Faza} \\
\midrule
A1 + topologia $\mathbb{Z}_2$
  & $\Rightarrow$
  & kink topologiczny, $n\in\mathbb{Z}$ (D.1a)
  & \statuslabel{Twierdzenie}
  & IV.A \\
D.1a + ZS1
  & $\Rightarrow$
  & spin $1/2$, statystyki Fermiego (D.1b--c)
  & \statuslabel{Propozycja}
  & IV.A \\
A1 + $a_\Gamma$, WKB
  & $\Rightarrow$
  & trzy pokolenia $N_{\mathrm{gen}}=3$
  & \statuslabel{Propozycja}
  & IV.B \\
$N_{\mathrm{gen}}=3$ + warunek bif.\ TGP
  & $\Rightarrow$
  & $Q_{\mathrm{Koide}} = 2/3$ (odtwarzanie)
  & \statuslabel{Propozycja}
  & IV.B \\
$Q=2/3$ + $r_{21}^{\mathrm{PDG}}$
  & $\Rightarrow$
  & $r_{31}$ (predykcja $\pm 0.016\%$); D.3
  & \statuslabel{Propozycja}
  & IV.B \\
OP-3 (otwarty)
  & $\Rightarrow$
  & $\alpha_K$ z~$a_\Gamma$ $\Rightarrow$ N$_\mathrm{param}=2$
  & \statuslabel{Program}
  & IV.C \\
\bottomrule
\end{tabular}
\end{center}

\begin{remark}[G\l{}\'owna hipoteza pomostowa]\label{rem:bridge-summary}
Centralnym \textbf{za\l{}o\.zeniem niewyprowadzalnym} pozosta\l{}ym w~TGP jest krok 3
mostu substrat~$\to$~metryka (uw.~\ref{rem:metric-bridge}):
\begin{equation*}
  g_{00} = -\frac{\Phi_0}{\Phi}
  \quad\text{z warunku}\quad
  c_{\mathrm{lok}}\cdot c_{\mathrm{graw}} = c_0^2.
\end{equation*}
Wszystkie \l{}a\'ncuchy I--IV s\k{a} poprawne pod warunkiem akceptacji
tego kroku.
Bez niego TGP dostarcza dynamik\k{e} pola $\Phi$, ale nie geometri\k{e}
Lorentzowsk\k{a}.
\end{remark}
